\begin{proof}
\begin{enumerate}[label=\arabic*.]
\item We use the following matrices on the marked dual lattice:
\[
A=\begin{pmatrix}
1&0&0&1&-1&1\\0&1&-1&-1&1&0\\0&1&0&-1&0&0\\
0&2&-1&-1&0&1\\0&2&0&-1&0&0\\0&1&1&-1&0&0
\end{pmatrix},\quad
B=\begin{pmatrix}
1&0&0&0&0&0\\0&0&-1&0&1&0\\0&0&-1&0&0&1\\
0&0&-1&0&0&0\\0&1&-1&0&0&0\\0&0&-1&1&0&0
\end{pmatrix}.
\]
Put $M_1=A^t$, $M_2=B^t$, $M_0=(M_1M_2)^{-1}$ and
\[
 v_1=(1,3,-2,-1,2,0)^t,\qquad v_2=(-1,0,0,0,0,0)^t.
\]
Theorem~\ref{night:P1.1} gives their orders, the primitive cusp
marking, the coinvariant homomorphism $\ell(x)=x_0$, and freeness of
every nonidentity power of the specified affine actions. Consequently
\[
 (T^6\times\Delta_u)/
 \langle(x,u)\mapsto(M_ix+v_i/m_i,\zeta_{m_i}u)\rangle,
 \qquad (m_1,m_2)=(5,4),
\]
are smooth oriented finite caps. On the universal angular cover
$u=r e^{2\pi ia}$ their actual marked collar maps are
\[
 \alpha_i(x,u)=x-av_i,
 \qquad
 \alpha_i(M_ix+v_i/m_i,\zeta_{m_i}u)=M_i\alpha_i(x,u).
\]
At the cusp we take the actual index-one Kuhn cap of
Theorem~\ref{night:P1.9}, in the marking of
Proposition~\ref{maps:src:F.2.1}. The three displayed boundary bundles
therefore glue to the marked torus bundle over the intervening pair
of pants, with zero additional clutch translation. Every boundary
is paired. The result is a closed connected smooth eight-manifold,
oriented by the cusp complex structure and the continued fibre/base
orientation.

\item The small cusp cap retracts onto its central fibre $C$ through
the actual mapping-cylinder homeomorphism
\cite[Theorem 4.1.1 and Lemmas 4.1.2--4.1.3, pp.~480--482]{ffAchingerOgus20}.
Indeed its divisorial log chart is
$\mathbb N\to\mathbb N^r$, $1\mapsto(1,\ldots,1)$; it is log smooth,
saturated and vertical, and the cap is proper by
Theorem~\ref{night:P1.9}. In the resulting cylinder the homotopy
$[x,r]\mapsto[x,(1-t)r]$ fixes $C$ and retracts onto it.

Before the period quotient, the components of $C$ and their
nonempty intersections are smooth complete toric strata. Each is
simply connected: a loop moves into its dense torus by general
position, and the one-parameter loops along the rays of a smooth
maximal cone bound the coordinate discs; those rays are an integral
basis. Van Kampen sends component-crossing paths to edges and triple
intersections to triangles of the Kuhn nerve in $\mathbb R^3$.
Thus the covering core has trivial fundamental group. Its period
quotient has deck group $\mathbb Z^3$, so $\pi_1(C)=\mathbb Z^3$.
The nearby-torus map is the quotient of $\Lambda$ by the saturated
image of $M_0-I$, by the actual $H_1$ collapse in
Theorem~\ref{night:P1.11}. It is onto, and the cusp meridian dies
because its zero section extends.

We also record the invariant normal closure needed below. Let
$K$ be the smallest $\langle M_1,M_2\rangle$-invariant subgroup
containing $(M_0-I)\Lambda$, and let
$J=(M_1-I)\Lambda+(M_2-I)\Lambda=\ker\ell$.
The equality $\ell M_i=\ell$ gives $K\subseteq\ker\ell$.
On the actual abelian quotient $\Lambda/K$, the induced maps $a,b$
satisfy $ab=1$, $a^5=1$ and $b^4=1$. Thus $a^4=a^5=1$, so
$a=b=1$ and $J\subseteq K$. Consequently $K=J=\ker\ell$ as
integral subgroups; the quotient argument uses no saturation.


Let $a_i$ be the two finite peripheral lifts. The universal covering
deck transformation of the $i$th cap has $m_i$th power equal to
translation by $v_i$, since $M_iv_i=v_i$. The actual collar map in
Step~1 sends $a_i$ to that transformation. Van Kampen therefore gives
\[
 a_i\lambda a_i^{-1}=M_i\lambda,
 \qquad a_1^5=v_1,\quad a_2^4=v_2,\quad a_1a_2=1.
\]
Write $a_1=t$ and $a_2=t^{-1}$. Then
$t=t^5t^{-4}=v_1+v_2$ lies in the abelian fibre subgroup. Conjugation
is consequently trivial, and that subgroup is a quotient of
$\Lambda_{\langle M_1,M_2\rangle}=\mathbb Z$, with coordinate $\ell$.
Now $\ell(v_1+v_2)=0$, so $t=0$, while $t^5=v_1$ kills its generator
$\ell(v_1)=1$. Every generator vanishes, proving $\pi_1=1$.
Equivalently, after this abelian reduction, the relations on
$(\ell,a_1,a_2)$ have matrix
\[
 \begin{pmatrix}-1&5&0\\1&0&4\\0&1&1\end{pmatrix},
 \qquad\det=-1.
\]

\item We calculate rational cohomology directly from the marked
restriction maps. For $1\leq q\leq4$ put
\[
 L_q=\bigwedge^q\Lambda^\vee\otimes\mathbb Q,
 \quad A_q=\bigwedge^qA,\quad B_q=\bigwedge^qB,
 \quad C_q=(B_qA_q)^{-1}.
\]
The finite restriction images are
$J_{1,q}=\ker(A_q-I)$ and $J_{2,q}=\ker(B_q-I)$: transfer followed
by restriction is multiplication by the finite covering degree,
which is invertible over $\mathbb Q$. The marked collapse in
Theorem~\ref{night:P1.11} gives
$J_{0,q}=\ker(C_q-I)$ at the cusp, with one additional rational
kernel in degree four. Proper base change identifies these with the
actual stalk maps
\cite[Tag 09V5, Lemma 20.18.1; chapter PDF p.~30]{StacksProject}.

Cutting the sphere along the two based peripheral arcs gives the
following cellular complex for the sheaf of restriction images:
\[
 J_{1,q}\oplus J_{2,q}\oplus J_{0,q}
 \xrightarrow{d_0}L_q^2\xrightarrow{\Delta_q}L_q,
\]
\[
 d_0(j_1,j_2,j_0)=(j_1-j_0,j_2-A_qj_0),
 \qquad
 \Delta_q(x,y)=(A_q-I)x+(B_q-I)y.
\]
The cusp value is transported across the first arc by $A_q$ before
the second subtraction. The relative disc boundary is $\Delta_q$;
their composition is zero because $B_qA_qj_0=j_0$.
The exterior matrices are their minors in lexicographically ordered
exterior bases. Rational elimination gives
\[
\begin{array}{c|r|rrr|rr}
q&\dim L_q&\dim J_{1,q}&\dim J_{2,q}&\dim J_{0,q}
 &\operatorname{rank}d_0&\operatorname{rank}\Delta_q\\\hline
1&6&2&2&3&6&6\\
2&15&3&3&9&15&15\\
3&20&4&4&9&17&20\\
4&15&3&3&9&15&15
\end{array}
\]
Here are explicit nonzero minors for the two rank calculations.
For each $J_{i,q}$ use the rational nullspace basis obtained from
reduced row echelon form: order the free columns increasingly, set
one free coordinate to one and the others to zero, and solve the
pivot coordinates. Number rows and columns from zero.
For $\Delta_q$ use all rows and the following columns:
\[
\begin{array}{c|l|r}
q&\text{columns}&\det\\\hline
1&1,2,3,5,7,8&2\\
2&0,1,2,4,\ldots,12,16,20,24&-6\\
3&0,\ldots,7,10,\ldots,17,20,24,30,36&-6\\
4&0,\ldots,7,10,11,12,13,15,21,25&6
\end{array}
\]
For $d_0$, use columns $(0,1,2,3,5,6)$ when $q=1$ and all
columns otherwise, with the following row sets:
\[
\begin{array}{c|l|r}
q&\text{rows}&\det\\\hline
1&0,1,3,4,6,7&1\\
2&0,\ldots,8,10,12,13,15,20,21&-2/3\\
3&0,\ldots,10,13,16,20,21,30,31&-4/3\\
4&0,\ldots,11,15,16,25&1/3
\end{array}
\]
All are determinants of the matrices specified above. They give
the stated full row ranks for $\Delta_q$ and full column ranks for
$d_0$ in degrees two through four. In degree one the common fixed
covector $e_0$ gives the nonzero kernel vector $(e_0,e_0,e_0)$ of
$d_0$, so its rank is at most six; the displayed minor gives equality.
The finite invariant dimensions count exterior products of the
specified eigenvalues with product one; the cusp dimensions are
those of Theorem~\ref{night:P1.11}.

These ranks determine the cohomology: $H^0$ of the image sheaf is
$\mathbb Q$ in coefficient degree one, $H^1$ is $\mathbb Q^3$ in
coefficient degree three, and every other entry in these four
coefficient degrees is zero. The cusp skyscraper kernel adds
$\mathbb Q$ to $H^0(R^4\pi_*\mathbb Q)$ and no higher cohomology.
Connected fibres give $R^0\pi_*\mathbb Q=\mathbb Q$.

In the Leray sequence
\cite[Tag 01F2, Lemma 20.13.4; chapter PDF p.~22]{StacksProject},
the only remaining low-degree differential is
$d_2^{0,1}:\mathbb Q\to\mathbb Q$. Its kernel is $H^1(M;\mathbb Q)=0$
by Step~2, so it is an isomorphism. Hence
$b_1=b_2=b_3=0$ and $b_4=3+1=4$. Poincar\'e duality determines all
remaining Betti numbers and gives \eqref{maps:eq:F.5}. Also, the
smooth and finite quotient fibres have Euler number zero, while the
Kuhn cusp has Euler number six by Theorem~\ref{night:P1.9}. Additivity
over the base strata therefore gives $\chi(M)=6$ directly.

\item Each finite monodromy admits an invariant real linear complex
structure: its $1$ and $-1$ real eigenspaces have even dimension, and
each remaining conjugate eigenpair gives a real two-plane. Choose
a complex structure on each summand. Reversing it on one fixed
two-plane selects the required orientation. These structures commute
with $M_i$, so their direct sums with the root-disc structure descend
to the actual affine caps. The cusp cap already has its complex
structure. Their vertical complex structures on the boundary collars
are translation-invariant in the marked torus fibres.

The space $\mathcal J=\operatorname{GL}_6^+(\mathbb R)/
\operatorname{GL}_3(\mathbb C)$ of allowable vertical structures
retracts to $\operatorname{SO}(6)/\operatorname{U}(3)$ by the explicit
conjugation
\[
 J\longmapsto P_tJP_t^{-1},
 \qquad P_t=\bigl((1-t)I+t(I+J^tJ)\bigr)^{1/2}.
\]
At $t=1$ the structure is orthogonal, since $J$ preserves the metric
$I+J^tJ$; an orthogonal $J$ is fixed throughout this homotopy.
The inclusion $\operatorname{U}(3)\to\operatorname{SO}(6)$ sends
$\operatorname{diag}(e^{i\theta},1,1)$ to one-plane rotation and
induces the surjection $\mathbb Z\to\mathbb Z/2$ on fundamental
groups. Indeed the determinant fibration and
$\operatorname{SU}(2)=S^3\to\operatorname{SU}(3)\to S^5$ give the
first group; the quaternion cover of $\operatorname{SO}(3)$ and
the sphere fibrations for $\operatorname{SO}(n)$ give the second,
with a full rotation lifting from $1$ to $-1$. The homogeneous-space
exact sequence therefore gives $\pi_0\mathcal J=\pi_1\mathcal J=0$.

Over the two-dimensional pair of pants the prescribed boundary
sections of the associated $\mathcal J$-bundle extend over relative
one-cells by connectedness and over relative two-cells by simple
connectivity. Smooth approximation relative to the collars gives
a smooth vertical structure. Horizontal splittings form an affine
space, so a partition of unity extends the prescribed complex
splittings from those collars. Combining the vertical structure with
the pulled-back base structure and gluing the caps gives the required
almost complex structure on $M$.

\item We have $c_1=c_3=0$ rationally because $H^2(M;\mathbb Q)$ and
$H^6(M;\mathbb Q)$ vanish, and $c_4[M]=\chi(M)=6$. The characteristic
number formulas in Theorem~\ref{maps:src:F.1.2}, with their source
\cite[equations (1.2)--(1.4), pp.~3--4]{ffEsoleKang}, give
\[
 \operatorname{Td}[M]=\frac{3c_2^2[M]-6}{720},
 \qquad
 \sigma(M)=\frac{3c_2^2[M]+84}{45}
           =16\operatorname{Td}[M]+2.
\]
The Todd genus of this almost complex manifold is an integer
\cite[Theorem 24.5.4 and corollary, p.~184]{HirzebruchTopologicalMethods}.
Since $|\sigma(M)|\leq b_4=4$, the displayed integer relation forces
$\operatorname{Td}[M]=0$ and $\sigma(M)=2$. Thus the real middle
pairing has positive and negative indices $3$ and $1$.

Integral Poincar\'e duality makes the middle pairing modulo torsion
unimodular. An even unimodular lattice has signature divisible by
eight, so this lattice is odd. The classification of indefinite odd
unimodular lattices
\cite[Theorem 2.2.1, p.~61]{ffPetersSterk} gives its integral isometry
class $3\langle1\rangle\oplus\langle-1\rangle$.
\end{enumerate}
\end{proof}
